\documentclass[11pt]{article}
\usepackage{amsmath,amssymb,amsthm,hyperref}
\newtheorem{lemma}{Lemma}
\newtheorem{corollary}{Corollary}
\newcommand{\E}{\mathbb E}
\newcommand{\Var}{\operatorname{Var}}
\title{Bias and variance of a success-weighted subset}
\author{Asymptotic research note; independently reviewed}
\date{30 September 2026}
\begin{document}\maketitle
This finite algebraic lemma controls the averaging error in a
forced-success polarization argument. It is independent of the
comparison-moment identities.

\begin{lemma}\label{main}
Let $p_1,\ldots,p_m\in(0,1]$, $1\le\ell<m$, $h=m-\ell$, and
choose an $\ell$-element subset $B$ with probability
\[
 \Pr(B)=\frac{\prod_{j\in B}p_j}{e_\ell(p)}.
\]
For any $a>0$, put $x_j=a(1-p_j)$, $\bar x=m^{-1}\sum_jx_j$,
$x_R=h^{-1}\sum_{j\notin B}x_j$, and $\Delta=x_R-\bar x$.
Then
\[
 \E\Delta=\frac{a}{mh}\sum_{i<j}(p_i-p_j)^2
       \frac{e_{\ell-1}(p_{-[i,j]})}{e_\ell(p)}\ge0.       \tag{1}
\]
If $b=\min_i p_i$ and $v=\sum_i(p_i-\bar p)^2$, then
\[
 0\le\E\Delta\le\frac{a v}{h b},\qquad
 \Var\Delta\le\frac{a^2}{h^2}\sum_i(1-p_i)^2.             \tag{2}
\]
In particular, writing $\delta=\max_i(1-p_i)$,
\[
 \E\Delta^2\le\frac{a^2m\delta^2}{h^2}
                 +\frac{a^2m^2\delta^4}{h^2b^2}.          \tag{3}
\]
\end{lemma}
\begin{proof}
Let $I_i=\boldsymbol1_{i\in B}$ and $\pi_i=\E I_i$.
The two inclusion probabilities have difference
\[
 \pi_i-\pi_j=(p_i-p_j)
       \frac{e_{\ell-1}(p_{-[i,j]})}{e_\ell(p)}.
\]
Also $\sum_i\pi_i=\ell$, and therefore
\[
 \sum_i x_i(\pi_i-\ell/m)
   =\frac1m\sum_{i<j}(x_i-x_j)(\pi_i-\pi_j).
\]
Since $\Delta=-(\sum_i x_i I_i-\ell\bar x)/h$ and
$x_i-x_j=-a(p_i-p_j)$, this proves (1).
The ratio in (1) is at most $1/p_i\le1/b$, since
$e_\ell(p)\ge p_i e_{\ell-1}(p_{-[i,j]})$.
Using $\sum_{i<j}(p_i-p_j)^2=mv$ proves its upper bound.

For the variance, pairwise inclusion covariances are nonpositive.
Here is a direct verification, requiring no general negative
association theorem. With $E_t=e_{\ell-t}(p_{-[i,j]})$ and
$Z=e_\ell(p)=E_0+(p_i+p_j)E_1+p_ip_jE_2$, one has
\[
 \operatorname{Cov}(I_i,I_j)
   =\frac{p_ip_j(E_0E_2-E_1^2)}{Z^2}\le0.
\]
The final inequality is Newton's elementary-symmetric inequality;
indices outside the available range are interpreted as zero.
Because $x_i\ge0$, it follows that
\[
 \Var\Bigl(\sum_i x_i I_i\Bigr)
 \le\sum_i x_i^2\Var(I_i)\le\sum_i x_i^2.
\]
This proves the variance bound in (2). Finally $v\le m\delta^2$
and $\sum_i(1-p_i)^2\le m\delta^2$, giving (3).
\end{proof}

\begin{corollary}[Application to endpoint comparison rows]\label{apply}
Assume the comparison-tensor hypotheses of
\path{finite_endpoint_odds_balance.tex}, $m\ge256$, and
$s=\sum_i(1-p_i)\le m/256$. Let $1\le\ell\le m/2$ and take
$a=\ell$ in Lemma~\ref{main}. There is an absolute constant $C$ such
that
\[
 0\le\E\Delta\le\frac{C(s+1)^2}{m},\qquad
 \E\Delta^2\le C\left(\frac{(s+1)^2}{m}
                          +\frac{(s+1)^4}{m^2}\right).    \tag{4}
\]
For every twice continuously differentiable function $P$ on
$[0,\ell s/h]$,
\[
 \left|\E P(x_R)-P(\bar x)\right|
 \le |P'(\bar x)|\frac{C(s+1)^2}{m}
 +\frac C2\sup_{0\le x\le\ell s/h}|P''(x)|
       \left(\frac{(s+1)^2}{m}+\frac{(s+1)^4}{m^2}\right).
 \tag{5}
\]
\end{corollary}
\begin{proof}
The cited finite balance theorem gives $p_i>0$ and
$r_i\le128(s+1)/m\le1$, hence $b\ge1/2$ and
$\delta\le128(s+1)/m$. Since $\ell\le h$ and $h\ge m/2$,
(2)--(3) imply (4). Both $x_R$ and $\bar x$ belong to the stated
interval, so Taylor's theorem and (4) give (5).
\end{proof}


\begin{corollary}[A refinement that vanishes at zero deficit]\label{refined}
In the setting of Corollary~\ref{apply},
\[
 0\le\E\Delta\le\frac{Cs(s+1)}m,\qquad
 \E\Delta^2\le C\left(\frac{s(s+1)}m
                         +\frac{s^2(s+1)^2}{m^2}\right). \tag{6}
\]
\end{corollary}
\begin{proof}
In (2), improve both quadratic sums by
\[
 v\le\sum_i(1-p_i)^2\le\delta s
                \le\frac{128s(s+1)}m.
\]
The same bounds on $\ell/h$ and $b$ then prove (6).
\end{proof}

\begin{lemma}[Taylor's remainder with a singular derivative envelope]\label{singular}
Let $x>0$, $v\ge0$, and let $P$ be twice continuously differentiable
on the interval with endpoints $x,v$. Suppose
\[
 |P''(z)|\le M e^{\gamma z}/z\qquad(z>0)
\]
there, with $\gamma\ge0$. Then
\[
 |P(v)-P(x)-P'(x)(v-x)|
       \le M e^{\gamma\max(x,v)}(v-x)^2/x.         \tag{7}
\]
\end{lemma}
\begin{proof}
Set $\Delta=v-x$. The integral form of Taylor's theorem uses
$\Delta^2\int_0^1(1-t)P''(x+t\Delta)\,dt$. If $\Delta\ge0$,
$(1-t)/(x+t\Delta)\le1/x$. If $\Delta<0$, then
$x+t\Delta=(1-t)x+tv\ge(1-t)x$, giving the same inequality for
$t<1$. Integration proves (7), including $v=0$; the apparent
singularity at the final integration endpoint is canceled by $1-t$.
\end{proof}

\begin{corollary}[The success-weighted mixture costs order $d/m$]\label{sharp}
Retain the comparison-tensor assumptions, take
$\ell=\lfloor m/4\rfloor$, $\alpha=\ell/m$, $x=\alpha s$,
and $\theta=1/16$. Suppose a polynomial $P$ satisfies, for $z>0$,
\[
 |P'(z)|\le C_0 B\sqrt{d/z}\,e^{\theta z},\qquad
 |P''(z)|\le C_0 B(d/z)e^{\theta z},               \tag{8}
\]
where $d\ge1$, $B\ge0$. On the rows $s\le m/256$,
\[
 (\ell+1)\E_w\left[
 \overline A\left|\E_{B\,\mathrm{weighted}}P(x_R)-P(x)\right|
 \boldsymbol1_{s\le m/256}\right]\le \frac{C B d}{m},       \tag{9}
\]
where the constant depends only on $C_0$.
\end{corollary}
\begin{proof}
If $s=0$, every $x_j$ and $\Delta$ is zero, so the error vanishes.
For $s>0$, use the mean term from (8) and the remainder bound (7).
Since $0\le x,x_R\le\ell s/h$ and $x=\alpha s$, (6) gives
\begin{align*}
 |\E P(x_R)-P(x)|
 &\le C B\sqrt d\,e^{\theta x}
              \frac{\sqrt s(s+1)}m\\
 &\quad+C Bd\,e^{\theta\ell s/h}
         \left(\frac{s+1}m+\frac{s(s+1)^2}{m^2}\right).
\end{align*}
Maclaurin's inequality gives $\overline A\le e^{-\alpha s}$.
The fixed choice of $\theta$ leaves a strictly negative exponential
factor after multiplication. The row-mass estimate in
\path{finite_endpoint_odds_balance.tex} implies
$\E_w (s+1)^k e^{-c s}\le C_{c,k}/m$ for each fixed $c>0,k$.
Apply it to the displayed terms, multiply by $\ell+1=O(m)$, and
use $\sqrt d\le d$. This proves (9).
\end{proof}

When averaging over uniformly chosen subsets, let
$\overline A=\binom m\ell^{-1}e_\ell(p)$ and
$A_B=\prod_{j\in B}p_j$. Then the left side of (5), multiplied by
$\overline A$, is exactly
\[
 \left|\E_{B\,\mathrm{uniform}}[A_BP(x_R)]
              -\overline A P(\bar x)\right|.
\]
Thus the first-order error is a bias of order $1/m$, not the order
$1/\sqrt m$ obtained by Cauchy--Schwarz. The note alone does not
assert a new bound on the number of dice faces; it supplies one
component of a proposed structured polynomial argument.
\end{document}
